%% exemplo-relatorio-externo.tex
%% Documento mínimo de teste — linha Relatório A4 CRP-SP, variante
%% `externo`: divulgação, com a camada de decoração.
%%
%% Exercita o que o `interno` não tem:
%%   - \relCapaArte + \relCapa, capa full-bleed SEM pdfpages;
%%   - \relDecoracao, faixas no topo e no pé por \AddToShipoutPictureBG;
%%   - a marcação de ARTEFATO das três, que é o ponto: nenhuma delas
%%     pode aparecer na árvore de tags como <Figure>.
%% E os créditos da relatorio.sty §12, que valem nas duas variantes:
%%   - página de créditos com título oculto (\creditosoculto), grupos
%%     com fio, subgrupos, duas colunas e cargos alinhados;
%%   - ficha técnica com título visível (\creditos).
%% Conferir nos créditos: um Span por pessoa com registro (12 aqui),
%% e o H1 "Créditos institucionais" na árvore sem texto pintado.
%% Os nomes são fictícios.
%%
%% ⚠️ A arte é de TESTE, gerada por arte-teste/gerar.sh — não é do
%% designer e não é conteúdo de publicação. Quando a arte real chegar,
%% muda só o caminho nos dois comandos abaixo.
%%
%% ⚠️ Aspas Unicode reais (“ ”), nunca `` '': a Atkinson não tem o
%% glifo da crase, e esta linha compila com \tracinglostchars=3.
%%
%% Compilar 2x com: lualatex exemplo-relatorio-externo.tex
%% Conferir:
%%   ~/verapdf/verapdf --format text --flavour ua2 exemplo-relatorio-externo.pdf
%%   texlua $(kpsewhich --progname=texlua show-pdf-tags.lua) --tree \
%%     exemplo-relatorio-externo.pdf | grep -cE '─Figure \('   # tem de dar 0
%%
%% ⚠️ O grep tem de casar o NÓ de estrutura (`─Figure (`), não a palavra
%% solta: este documento escreve "Figure" na prosa e numa célula de
%% tabela, e um `grep -c Figure` devolve 3 com a árvore perfeitamente
%% limpa. O falso positivo já custou uma investigação.

\DocumentMetadata{
  lang        = pt-BR,
  pdfversion  = 2.0,
  pdfstandard = ua-2,
  tagging     = on,
}

\documentclass[externo]{crpsp-relatorio}

\hypersetup{
  pdftitle  = {Exemplo de Relatório Externo CRP-SP},
  pdfauthor = {Conselho Regional de Psicologia de São Paulo},
}

\title{Exemplo de Relatório Externo}
\author{Conselho Regional de Psicologia de São Paulo}
\date{Setembro de 2026}

% A camada do designer, declarada por arquivo. Faltando o arquivo, a
% faixa fica vazia com aviso e a capa cai na versão tipográfica — é o
% que permite adiantar a classe antes de a arte existir.
\relCapaArte{arte-teste/capa.pdf}
\relDecoracao{arte-teste/faixa-topo.pdf}{arte-teste/faixa-pe.pdf}

\begin{document}

\relCapa

% Créditos: a página não imprime título, mas o leitor de tela navega
% por ele. Os grupos, as quebras e a ordem são do documento; as linhas
% são o que o gerador do normativas-pipeline escreve.
\creditosoculto{Créditos institucionais}
\creditosgrupo{Plenário}
\creditossub{Diretoria}
\begin{CreditoInstitucional}
  \CargoNosCreditos[06/00001]{Pessoa de Teste Um}{presidenta}
  \CargoNosCreditos[06/00002]{Pessoa de Teste Dois, com um nome
    longo que quebra dentro da própria coluna}{vice-presidente}
  \CargoNosCreditos[06/00003]{Pessoa de Teste Três}{secretária}
\end{CreditoInstitucional}
\creditossub{Conselheiras/os}
\begin{NomesDuasColunas}
  \NomeNosCreditos{Pessoa de Teste Quatro}{06/00004}
  \NomeNosCreditos{Pessoa de Teste Cinco}{06/00005}[licenciada]
  \NomeNosCreditos{Pessoa de Teste Seis}{06/00006}
  \NomeNosCreditos{Pessoa de Teste Sete}{06/00007}
  \NomeNosCreditos{Pessoa de Teste Oito}{06/00008}
  \NomeNosCreditos{Pessoa de Teste Nove}{06/00009}
\end{NomesDuasColunas}
\creditosgrupo{Gerência de Teste}
\begin{CreditoInstitucional}
  \CargoNosCreditos[06/00010]{Pessoa de Teste Dez}{gerente}
  \CargoNosCreditos{Pessoa de Teste Onze}{coordenadora de Tecnologia
    da Informação e Comunicação}
  \CargoNosCreditos[06/00012]{Pessoa de Teste Doze}{coordenador de
    Comunicação}
\end{CreditoInstitucional}

\section{A decoração}

A capa acima e as faixas no topo e no pé desta página são arte de
teste. As três entram como \emph{artefato}: o leitor de tela não as
encontra na ordem de leitura, e a árvore de tags não traz nenhum
\texttt{Figure} vindo delas.

Segundo parágrafo, para conferir que o fluxo de texto passa por cima
da decoração sem colidir com ela — a mancha tem 35\,mm de margem em
todos os lados, e é essa faixa que a arte ocupa.

\subsection{O que muda em relação ao interno}

\begin{itemize}
  \item A variante \texttt{interno} não carrega \texttt{eso-pic} e não
        define capa nem faixas.
  \item A variante \texttt{externo} define as duas coisas, e ambas
        degradam para nada quando o arquivo de arte falta.
  \item As duas são mutuamente exclusivas: combinadas, a classe emite
        erro.
\end{itemize}

\section{O resto continua igual}

\begin{destaque}{Critério de aceitação}
  A decoração só está certa se o veraPDF der \textbf{PASS em ua2} e a
  árvore de tags não tiver nenhum \texttt{Figure} vindo da arte.
  Passar por zero erro fatal não basta: decoração na árvore de leitura
  é defeito silencioso.
\end{destaque}

\begin{tabular}{L{0.4\linewidth}L{0.45\linewidth}}
  \linhacab \celcab{Elemento} & \celcab{Marcação} \\
  Faixa do topo & Artefato \\
  Faixa do pé & Artefato \\
  \linhadestaque \textbf{Capa} & \textbf{Artefato} \\
  Gráfico no fluxo & \texttt{Figure} com \texttt{alt} \\
\end{tabular}

Parágrafo final, fechando o documento de teste.

\creditos{Ficha técnica}
\begin{CreditoInstitucional}
  \NomeNosCreditos{Pessoa de Teste Treze}{06/00013}
  Diagramação e revisão de acessibilidade.
\end{CreditoInstitucional}

\end{document}
